\documentclass[12pt, leqno]{article}
\usepackage[english,russian]{babel}
\usepackage{fontspec}
\usepackage[utf8]{inputenc}
\usepackage{epigraph}
\usepackage{tabularx}
\usepackage[leqno]{amsmath}
\setmainfont{Times New Roman}
\setlength\parindent{0pt}
\usepackage[colorlinks=true,urlcolor=blue]{hyperref}        
\usepackage{minted} 
\usepackage[xetex,final]{graphicx}
\usepackage{float}%"Плавающие" картинки
\usepackage{wrapfig}%Обтекание фигур (таблиц, картинок и прочег
\usepackage{epstopdf}
\usepackage{amssymb}
 
\begin{document}

\section*{Числа фибоначчи. Рекурсия}
\begin{flushleft}
\textbf{Вопросы рекурсии:} Рекурентные формулы и реализация
\end{flushleft}

Последовательность чисел Фибоначчи  $ F_{n} $ задаётся линейным рекуррентным соотношением: 

$ F_{0}=0,\quad F_{1}=1,\quad F_{n}=F_{n-1}+F_{n-2} $ где $ n \geqslant 2,\ n\in \mathbb {Z}$ 
 

На Западе эта последовательность была исследована Леонардо Пизанским, известным как Фибоначчи, в его труде «Книга абака» (1202) . Он рассматривает развитие идеализированной (биологически нереальной) популяции кроликов, где условия таковы: изначально дана новорождённая пара кроликов (самец и самка); со второго месяца после своего рождения кролики начинают спариваться и производить новую пару кроликов, причём уже каждый месяц; кролики никогда не умирают, — а в качестве искомого выдвигает количество пар кроликов через год:

\begin{enumerate}
    \item начале первого месяца есть только одна новорождённая пара (1);\\
    \item конце первого месяца по-прежнему только одна пара кроликов, но уже спарившаяся (1);\\
    \item конце второго месяца первая пара рождает новую пару и опять спаривается (2);\\
    \item конце третьего месяца первая пара рождает ещё одну новую пару и спаривается, вторая пара только спаривается (3).\\
    \item конце четвёртого месяца первая пара рождает ещё одну новую пару и спаривается, вторая пара рождает новую пару и спаривается, третья пара только спаривается (5).\\
\end{enumerate}

В конце n-го месяца количество пар кроликов будет равно количеству пар в предыдущем месяце плюс количеству новорождённых пар, которых будет столько же, сколько пар было два месяца назад, то есть $ F_{0}=0,\quad F_{1}=1,\quad F_{n}=F_{n-1}+F_{n-2} $. Возможно, эта задача также оказалась первой, моделирующей экспоненциальный рост популяции. 
 
 

\clearpage 
 
Дальше можно смотреть по ссылке:

\href{https://ru.wikipedia.org/wiki/Числа_Фибоначчи}{Числа Фибоначчи}


\clearpage 

\noindent Образец кода на лисп (sbcl):  
\begin{minted}[fontsize=\small]{lisp}

;;;; Programm for calculation Fibnacci sequences
;;;; https://www.jdoodle.com/execute-clisp-online

;;; FIB computes the the Fibonacci function in the traditional
;;; recursive way.


;; Fibonacci  calc function. 
(defun fib(n)
    (check-type n integer)
    ;; At this point we're sure we have an integer argument.
    ;; Now we can get down to some serious computation.
    (cond  ((= n 0) 0)
           ((= n 1) 1)
           (t ( +  (fib (- n 2))       ; The traditional formula
                   (fib (- n 1)) )))   ; is fib[n-1]+fib[n-2].
)

(defun cli/parameter()
;;Get command line argument" 
    (let (
        (args sb-ext:*posix-argv*))
        (car (cdr args)))
)

(defun makeDomain(num acc)
;;Make domain's interval from 0 to num " 
    (cond ((eq num 0) acc)
          (t (makeDomain (- num 1) (cons num acc))))
)

;; call fib  function for  interval
;; input   ./exec4.lisp  <n = 25>
;; main print
(defun main()
    (terpri)
    (format T "~a" (mapcar 'fib (makeDomain (parse-integer (cli/parameter)) nil)))
)

(main)

;;--------------------------

\end{minted}  


\noindent Образец кода на haskell:


\begin{minted}[fontsize=\small]{haskell}
module Math where

fib :: Int -> Int
fib 0 = 0
fib 1 = 1
fib n = fib (n-1) + fib (n-2)


-- The memoized version is much faster. Try memoized_fib 10000.
-- http://wiki.haskell.org/Memoization
-- http://lin-techdet.blogspot.com/2015/07/haskell.html

memoized_fib :: Int -> Integer
memoized_fib = (map fib [0 ..] !!)
   where fib 0 = 0
         fib 1 = 1
         fib n = memoized_fib (n-2) + memoized_fib (n-1)
\end{minted}  

Мain.hs:
\begin{minted}[fontsize=\small]{haskell}
module Main where

import System.Environment   
import Data.List  
import Math

{- https://stackoverflow.com/questions/1105765/generating-fibonacci-numbers-in-haskell -}

getIntArg :: IO Int
getIntArg = fmap (read . head) getArgs

main = do  
    args <- getIntArg              -- IO [Int]
    putStrLn "The arguments are:"  
    print (fib args)
\end{minted} 

\clearpage 
Ниже представлен результат для $ n = 20$


0, 1, 1, 2, 3, 5, 8, 13, 21, 34, 55, 89, 144, 233, 377, 610, 987, 1597, 2584, 4181, 6765, 10946, 17711

 

  

\end{document}